\section{Jejak Memori}

Hipotesis H2 menyatakan bahwa jejak memori CoDA-EDA tidak bertumbuh terhadap ukuran populasi virtual $(NP)$, berbeda dari algoritma yang menyimpan populasi eksplisit. 
Pengujian dilakukan pada $NP\in{25,50,100,200,400}$ menggunakan $MemPeak_{alloc}$ dari $tracemalloc$, dengan setiap konfigurasi dijalankan dalam proses terpisah. 
Hasil rerata lima dataset disajikan pada Tabel \ref{tab:skala-memori-np}.

\begin{longtable}{|>{\raggedright\arraybackslash}p{2cm}|>{\raggedright\arraybackslash}p{2cm}|>{\raggedright\arraybackslash}p{2cm}|>{\raggedright\arraybackslash}p{2cm}|>{\raggedright\arraybackslash}p{4cm}|}
\caption{Penskalaan Jejak Memori terhadap Ukuran Populasi NP (rerata lintas lima dataset)}
\label{tab:skala-memori-np} \\
\hline
\textbf{Algoritma} & \textbf{alloc NP=25 (KB)} & \textbf{alloc NP=400 (KB)} & \textbf{Rasio 400/25} & \textbf{Keterangan} \\
\hline
\endfirsthead

\multicolumn{5}{c}{{\bfseries \tablename\ \thetable{} -- lanjutan dari halaman sebelumnya}} \\
\hline
\textbf{Algoritma} & \textbf{alloc NP=25 (KB)} & \textbf{alloc NP=400 (KB)} & \textbf{Rasio 400/25} & \textbf{Keterangan} \\
\hline
\endhead

\hline
\multicolumn{5}{r}{Lanjut ke halaman berikutnya} \\
\endfoot

\hline
\endlastfoot

% ==================== DATA BARIS ====================
CoDA-EDA & 492 & 492 & 1{,}00$\times$ & DATAR --- memori tak bergantung NP ($O(D)$) \\
\hline
compact-GA & 169 & 103 & 0{,}77$\times$ & Datar (compact univariat) \\
\hline
compact-PSO & 172 & 106 & 0{,}76$\times$ & Datar (compact univariat) \\
\hline
iBOA & 190 & 134 & 0{,}82$\times$ & Datar (statistik inkremental) \\
\hline
ACO & 212 & 2.036 & 6{,}22$\times$ & Tumbuh --- populasi/feromon \\
\hline
BOA & 349 & 4.028 & 7{,}68$\times$ & Tumbuh --- populasi + jaringan Bayes \\
\hline
UBK & 387 & 5.169 & 9{,}20$\times$ & Tumbuh --- populasi swarm \\
\hline

\end{longtable}

Pola yang diperoleh konsisten dengan mekanisme penyimpanan masing-masing algoritma. 
Ketika $NP$ dinaikkan enam belas kali dari 25 menjadi 400, puncak alokasi CoDA-EDA tetap sekitar 492 KB dengan rasio $R_{NP}=1{,}00$. 
Hal ini terjadi karena $NP$ hanya menentukan granularitas pembaruan probabilistik, bukan jumlah kandidat yang disimpan. 
Sebaliknya, ACO, BOA, dan UBK menyimpan struktur yang berkaitan dengan populasi eksplisit sehingga puncak alokasinya meningkat masing-masing $6,22$, $7,68$, dan $9,20$ kali.

Keluarga \textit{compact} dan iBOA menunjukkan pola serupa dengan CoDA-EDA, dengan $R_{NP}$ masing-masing sekitar $0,76$–$0,77$ dan $0,82$. 
Nilai di bawah 1 mencerminkan variasi pengukuran, tetapi secara keseluruhan menunjukkan tidak adanya pertumbuhan sistematis terhadap $NP$. 
Dengan demikian, hasil empiris mendukung perbedaan antara algoritma tanpa populasi eksplisit dan algoritma yang kebutuhan memorinya bertambah ketika ukuran populasi diperbesar.

Sebagai bukti pelengkap terhadap ketergantungan memori pada dimensi, Gambar \ref{fig:jejak-memori} memetakan $MemPeak_{alloc}$ CoDA-EDA terhadap $D$ pada lima dataset. 
Regresi menghasilkan hubungan sekitar
$
MemPeak_{alloc}\approx0{,}80D+93\text{ KB},
$
dengan $R^2=1{,}000$\label{first:symbol-r2} dan korelasi Pearson $r=1{,}000$, sedangkan variasi memori terhadap $NP$ berada di bawah $0,05$ KB. 
Hasil ini konsisten dengan analisis \textit{by-construction} bahwa kebutuhan memori CoDA-EDA bertumbuh linear terhadap $D$ dan tidak bergantung pada $NP$. 
Namun, karena lima titik $D$ berasal dari dataset berbeda dan bukan dari \textit{sweep} dimensi terkontrol pada dataset yang sama, hasil tersebut diperlakukan sebagai validasi empiris pendukung, bukan pembuktian asimtotik tersendiri.

Secara absolut, CoDA-EDA menggunakan memori lebih besar daripada \textit{compact-GA/compact-PSO} ($103$-$172$ KB) dan iBOA ($134$–$190$ KB), dengan puncak sekitar $492$ KB akibat tambahan struktur dependensi. 
Namun, fokus H2 adalah perilaku penskalaan, bukan penggunaan memori terendah pada satu konfigurasi. 
CoDA-EDA mempertahankan pemodelan dependensi tanpa mengalami pertumbuhan memori terhadap $NP$, sedangkan algoritma berpopulasi eksplisit menunjukkan peningkatan $6,22$–$9,20$ kali pada rentang pengujian.

\begin{figure}[H]
    \centering
    \includegraphics[width=0.9\textwidth]{Gambar_Memori_vs_D_CoDA-EDA.png}
    \caption{Jejak Memori CoDA-EDA Linear Terhadap Dimensi Ruang Solusi D}
    \label{fig:jejak-memori}
\end{figure}

Dengan demikian, H2 didukung pada konfigurasi yang diuji. 
CoDA-EDA mempertahankan jejak memori yang stabil ketika $NP$ dinaikkan, sekaligus menunjukkan pola pertumbuhan terhadap $D$ yang konsisten dengan kompleksitas memori $O(D)$. 
Hasil ini menunjukkan bahwa pemodelan dependensi berorde terbatas dapat ditambahkan ke mekanisme \textit{compact} tanpa menghilangkan sifat utamanya, yaitu tidak menyimpan populasi virtual secara eksplisit.